\begin{ouverture}[courbe]{7}{Grandeurs composées et conversions}{Partie 2 · Grandeurs et fonctions}
{Durées · Vitesses · Débits · Aires · Cohérence des unités}
\objectif{Convertir les unités adaptées à une situation.}
\objectif{Calculer avec une vitesse, un débit ou une masse volumique.}
\objectif{Contrôler les dimensions d’un résultat.}
\prerequis{Proportionnalité et unités de longueur.}
\end{ouverture}
\pagestyle{otMain}\essentiel\fontsize{11}{15}\selectfont
\essSection{1}{Choisir et convertir les unités}
\begin{propriete}Pour les longueurs, $1\text{ m}=100\text{ cm}$. Pour les aires, $1\text{ m}^2=10000\text{ cm}^2$ ; pour les volumes, $1\text{ m}^3=10^6\text{ cm}^3$. Le facteur de conversion est mis au carré ou au cube.\end{propriete}
\begin{propriete}$1\text{ L}=1\text{ dm}^3$, $1\text{ mL}=1\text{ cm}^3$ et $1\text{ m}^3=1000\text{ L}$. Une tonne vaut 1000 kg. Une heure vaut 60 min, soit 3600 s.\end{propriete}
\begin{exemple}$2{,}4\text{ h}=2\text{ h}+0{,}4\times60\text{ min}=2\text{ h }24\text{ min}$. À l’inverse, 2 h 15 min valent $2+15/60=2{,}25$ h.\end{exemple}
\begin{attention}1 h 30 min n’est pas 1,30 h. De même, multiplier par 100 pour convertir une aire en m$^2$ vers des cm$^2$ est incorrect : le facteur est 10000.\end{attention}
\begin{savoirfaire}[Contrôler la cohérence]Nommer la grandeur et son unité à chaque étape. Une aire ne s’exprime pas en mètres, une durée ne s’exprime pas en kilomètres. Donner un ordre de grandeur avant de conclure.\end{savoirfaire}
\clearpage
\essSection{2}{Des rapports de grandeurs}
\begin{propriete}Une vitesse moyenne vaut $v=d/t$, avec $t>0$. On peut utiliser $d=vt$ et $t=d/v$ pour $v>0$. Une vitesse moyenne sur un trajet se calcule avec la distance totale et la durée totale.\end{propriete}
\begin{exemple}$72\text{ km/h}=72000/3600=\res{20\text{ m/s}}$. Pour convertir des km/h en m/s, diviser par 3,6 ; pour l’opération inverse, multiplier par 3,6.\end{exemple}
\begin{definition}Un débit volumique est un volume par unité de temps : $q=V/t$. Une masse volumique est une masse par unité de volume : $\rho=m/V$. Choisir les unités cohérentes avant de calculer.\end{definition}
\begin{savoirfaire}[Une chaîne de calculs]Pour remplir 600 L à un débit de 15 L/min, la durée est $600/15=40$ min. À débit constant, doubler le volume double la durée. Pour une consommation en kWh, multiplier une puissance en kW par une durée en heures.\end{savoirfaire}
\begin{attention}La moyenne de deux vitesses n’est pas généralement la vitesse moyenne d’un trajet. Les durées passées à chaque vitesse comptent.\end{attention}
\clearpage\rubrique{Découverte}
\begin{blocLivre}{ex-7-1}\begin{exercice}[1]{Des longueurs et des aires}
Convertir 3,8 km en mètres, 0,42 m$^2$ en cm$^2$ et 6500 cm$^2$ en m$^2$.
\end{exercice}\end{blocLivre}
\begin{blocLivre}{ex-7-2}\begin{exercice}[1]{Volumes et contenances}
Convertir 2,7 m$^3$ en litres et 450 mL en cm$^3$ puis en litres.
\end{exercice}\end{blocLivre}
\begin{blocLivre}{ex-7-3}\begin{exercice}[1]{Une durée décimale}
Écrire 1,75 h en heures et minutes ; écrire 2 h 42 min en heures décimales.
\end{exercice}\end{blocLivre}
\begin{blocLivre}{ex-7-4}\begin{exercice}[1]{Une vitesse moyenne}
Un trajet de 150 km est parcouru en 2 h 30 min. Calculer la vitesse moyenne.
\end{exercice}\end{blocLivre}
\begin{blocLivre}{ex-7-5}\begin{exercice}[1]{Un débit}
Un robinet fournit 84 L en 7 minutes, à débit constant. Donner son débit et le volume en 25 minutes.
\end{exercice}\end{blocLivre}

\clearpage\rubrique{Entraînement}
\begin{blocLivre}{ex-7-6}\begin{exercice}[2]{Convertir une vitesse}
Convertir 90 km/h en m/s et 12 m/s en km/h.
\end{exercice}\end{blocLivre}
\begin{blocLivre}{ex-7-7}\begin{exercice}[2]{Le temps de réaction}
Une voiture roule à 54 km/h. Le conducteur met 0,8 s avant de commencer à freiner. Quelle distance parcourt-elle pendant cette réaction ?
\end{exercice}\end{blocLivre}
\begin{blocLivre}{ex-7-8}\begin{exercice}[2]{Une masse volumique}
Un matériau a une masse volumique de 2,5 g/cm$^3$. Un bloc occupe 360 cm$^3$. Calculer sa masse en grammes puis en kilogrammes.
\end{exercice}\end{blocLivre}
\begin{blocLivre}{ex-7-9}\begin{exercice}[2]{Une consommation électrique}
Un appareil de puissance constante 1,2 kW fonctionne pendant 45 minutes. Calculer l’énergie en kWh. Le kWh coûte 0,24 € dans le modèle proposé. Calculer le coût.
\end{exercice}\end{blocLivre}
\begin{blocLivre}{ex-7-10}\begin{exercice}[2]{Débit en deux unités}
Convertir 0,35 m$^3$/s en L/min. Quel volume passe en 4 minutes ?
\end{exercice}\end{blocLivre}

\clearpage\rubrique{Problèmes}
\begin{blocLivre}{ex-7-11}\begin{exercice}[3]{Aller et retour}
Un cycliste parcourt 12 km à 18 km/h, puis revient sur 12 km à 12 km/h. Calculer sa vitesse moyenne sur l’aller-retour.
\end{exercice}\end{blocLivre}
\begin{blocLivre}{ex-7-12}\begin{exercice}[3]{Deux robinets}
Un premier robinet remplit seul une cuve en 30 min, un second en 45 min. Ils fonctionnent ensemble avec des débits constants et indépendants. Quelle durée pour remplir la cuve ?
\end{exercice}\end{blocLivre}
\begin{blocLivre}{ex-7-13}\begin{exercice}[3]{Choisir assez de peinture}
Un mur rectangulaire mesure 5 m sur 2,6 m. Une ouverture de 2 m$^2$ n’est pas peinte. Deux couches sont prévues. Un litre couvre 8 m$^2$ par couche. Les pots contiennent 2,5 L. Combien de pots acheter ?
\end{exercice}\end{blocLivre}
\begin{blocLivre}{ex-7-14}\begin{exercice}[3]{Une fuite discrète}
Une fuite laisse tomber 2 gouttes par seconde. On modélise une goutte par 0,05 mL. Quel volume d’eau est perdu en une journée ? En 30 jours ?
\end{exercice}\end{blocLivre}
\begin{blocLivre}{ex-7-15}\begin{exercice}[3]{Une livraison de gravier}
Une allée de 12 m sur 1,5 m reçoit une couche uniforme de gravier de 6 cm. La masse volumique apparente est 1,6 tonne/m$^3$. Un camion peut emporter 1,5 tonne. Combien de voyages au minimum ?
\end{exercice}\end{blocLivre}

\clearpage
